% ---------------------------------------------------------------------------
% blog/src/2026-09-21-typhoon-hil-techday-2026.tex
% Typhoon HIL TechDay at the MIT Samberg Conference Center, 21--22 September
% 2026. Datacenter power, a very small hotel room, and the Great Dome.
% ---------------------------------------------------------------------------

\title{Typhoon HIL TechDay 2026}
\date{2026-09-21 to 2026-09-22}
\location{Boston, MA}
\tags{travel, life}

\begin{document}

\lead{I spent two days at the annual Typhoon HIL TechDay, held this year at
MIT, as a master's grad scouting for somewhere to do a PhD. It was datacenter
power by day and the smallest hotel room in Back Bay by night.}

I went to Boston with two agendas. The official one was the conference. The
unofficial one is that I have a master's degree now and no PhD program yet,
and MIT is, of course, a place of interest. Getting to do the first while
standing on the campus of the second was a fortunate way to spend a week.

\section{The conference}

The thread running through both days was how the market is getting ready to
deliver power to what Typhoon HIL's Dr. Nikola Celanovic calls ``Knowledge
Factories.'' Everyone else calls them datacenters. His name for them stuck
with me.

\img{/assets/tecgDAylinkedin.jpg}{A question from the floor at the Samberg
Conference Center, with the Boston skyline doing its best to distract.
Photo from Typhoon HIL on LinkedIn.}

The keynote I got the most out of came from Dr. Himanshu Prasad of Schneider
Electric. He covered a large portion of where the industry stands today on
how power is delivered inside a datacenter, and then laid out what the next
phases of solid-state transformer (SST) adoption look like. It is one thing
to read about SSTs in a paper and another to hear someone whose company has
to ship them explain the order in which the pieces will actually land.

The other keynote that stayed with me was Dr. Stuart Laval of Eaton, on how
model-based engineering has changed the way the industry works. The case he
made is simple to state and hard to argue with: put enough fidelity into the
model and the simulation, and you can wring out the controller completely
before it ever meets real hardware. Sitting in a room full of people who do
exactly that for a living made the point land harder than any slide could.

\img{/assets/sambergskylineview.jpg}{The view from the Samberg Conference
Center, across the Charles toward Beacon Hill, Longfellow Bridge on the left.
Hard to focus on a slide deck with this behind it.}

\sep

\section{Boston between sessions}

When I was not sitting in a session, I got to see a bit of Boston. Allow me
to overreact for a moment: the Charlesmark Hotel in Back Bay is currently my
favorite hotel in the world. It is unbelievably quaint. They gave me the
wrong room key when I checked in. The in-room entertainment was a ten-inch
television from the Cold War, I am sure of it. My view was of an alley where
at least three stabbings happened this month, probably. I absolutely love
it.

\img[50]{/assets/charlesmarkhotel.jpg}{The room at the Charlesmark. The
television is just out of frame, which is about the right amount of frame
for it.}

It is nestled in and cozy, near enough to people and things, and the moment
you leave the building the city is right there, no transition. It does not
hurt that there is a Blank Street coffee shop next door.

\img{/assets/blankstboston.jpg}{Blank Street, next door. Stained glass in a
coffee shop, because Boston.}

\img{/assets/bay.jpg}{Back Bay from the Charles at dusk, the Prudential lit
up on the left.}

\sep

\section{Paying homage at MIT}

I did, in fact, pay homage and visit the MIT campus. I was not able to meet a
professor on this trip, but I did reconnect with a colleague who was at
Illinois with me and is now doing his PhD at MIT. If he is reading this:
THANKS MAN! I really appreciate it.

\img{/assets/mitdome.jpg}{The Great Dome at night, MCMXVI over the door.}

Nothing gets me amped up like walking into a new lab and asking questions
about every project that is going on in it. Seeing that lab alone made the
trip worth it.

\end{document}
